\subsection{中心函数的 Fourier 变换}

对于所有 \(u \in \hat{\mathrm{G}}\)，记 \(\chi_u\) 为 \(u\) 的特征标；于是

\[
\chi_ {u} (g) = \operatorname{Tr} (u (g)), \quad (g \in \mathrm{G}).\tag{16}
\]

回忆 Spectral Theory 中的公式

\[
\chi_ {u} * \chi_ {v} = 0 (u, v \in \hat {\mathrm{G}}, u \neq v)\tag{17}
\]

\[
\chi_ {u} * \chi_ {u} = \frac {1}{d (u)} \chi_ {u} \quad (u \in \hat {\mathrm{G}}).\tag{18}
\]

对于所有 \(u \in \hat{G}\)，记 \(\varepsilon_{u}\) 为 \(E_{u}\) 的恒等映射。回忆 (§7, no. 4)，\(\mathrm{ZL}^{2}(\mathrm{G})\) 表示 \(\mathrm{L}^{2}(\mathrm{G})\) 的子空间，它由中心函数 f 的等价类组成；这里中心函数是指满足 \(f \circ \operatorname{Int}s = f\) 对所有 \(s \in G\) 成立的函数，或者等价地，满足 \(\gamma(s)f = \delta(s^{-1})f\) 对所有 \(s \in G\) 成立的函数。

命题 4. 设 \(f \in \mathrm{L}^2(\mathrm{G})\)。则 \(f\) 是中心的，当且仅当 \(u(f)\) 对所有 \(u \in \hat{\mathrm{G}}\) 都是数乘变换。在此情形下

\[
u (f) = \frac {\varepsilon_ {u}}{d (u)} \int_ {\mathrm{G}} f (g) \chi_ {u} (g) d g.\tag{19}
\]

由命题 1 (no. 1)，\(f\) 为中心的意思是 \(u(\gamma(s)f) = u(\delta(s^{-1})f)\) 对所有 \(s \in \mathbf{G}\) 和所有 \(u \in \hat{\mathbf{G}}\) 成立；但这也可写成 \(u(s)u(f) = u(f)u(s)\) 对所有 \(s \in \mathbf{G}\) 和所有 \(u \in \hat{\mathbf{G}}\) 成立（公式 (12) 和 (13)），故命题 4 的第一条断言成立（Schur 引理）。若 \(u(f)\) 是数乘变换，则 \(u(f) = \lambda_u\varepsilon_u\)，其中

\[
\lambda_ {u} = \frac {1}{d (u)} \operatorname{Tr} (u (f)) = \frac {1}{d (u)} \int_ {\mathrm{G}} f (g) \operatorname{Tr} (u (g)) d g = \frac {1}{d (u)} \int_ {\mathrm{G}} f (g) \chi_ {u} (g) d g.
\]

因此，对于 \(f \in \mathrm{ZL}^2(\mathbf{G})\)，有

\[
u (f) = \langle \overline {{\chi}} _ {u} | f \rangle \frac {\varepsilon_ {u}}{d (u)} \rangle ,\tag{20}
\]

因此

\[
\overline {{{{\mathcal {F}}}}} (f) = \left(\langle \overline {{{{\chi}}}} _ {u} | f \rangle \frac {\varepsilon_ {u}}{d (u)}\right) _ {u \in \hat {\mathrm{G}}}\tag{21}
\]

其中

\[
\| \overline {{\mathcal {F}}} (f) \| _ {2} ^ {2} = \sum_ {u} \left| \left| \langle \overline {{\chi}} _ {u} | f \rangle \frac {\varepsilon_ {u}}{d (u)} \right| \right| _ {2} ^ {2} = \sum_ {u} | \langle \overline {{\chi}} _ {u} | f \rangle | ^ {2}.
\]

反之，若 \(\varphi\) 是 \(\hat{\mathbf{G}}\) 上的平方可积复函数，则元素 \(\left(\frac{\varphi(u)}{d(u)}\varepsilon_u\right)_{u\in \hat{\mathbf{G}}}\) 属于 \(\mathrm{F}(\hat{\mathrm{G}})\) 及 \(\mathrm{L}^2 (\hat{\mathrm{G}})\)，并且由公式 (9) 有

\[
\left(\mathcal {F} _ {u} \left(\frac {\varphi (u)}{d (u)} \varepsilon_ {u}\right)\right) (g) = d (u) \mathrm{Tr} \left(\frac {\varphi (u)}{d (u)} \varepsilon_ {u} u (g) ^ {- 1}\right) = \varphi (u) \overline {{\chi}} _ {u} (g),
\]

因此

\[
\mathcal {F} \left(\left(\frac {\varphi (u)}{d (u)} \varepsilon_ {u}\right)\right) = \sum_ {u \in \hat {\mathrm{G}}} \varphi (u) \overline {{\chi}} _ {u}.\tag{22}
\]

特别地，注意公式 (20) 和 (21) 给出，对于 u、v 属于 \(\hat{G}\)，

\[
u (\overline {{\chi}} _ {v}) = 0 \quad \text { if } u \neq v,\tag{23}
\]

\[
u (\overline {{\chi}} _ {u}) = \frac {\varepsilon_ {u}}{d (u)} \in \mathrm{End} (\mathrm{E} _ {u}),\tag{24}
\]

\[
\overline {{\mathcal {F}}} (\chi_ {u}) = \frac {\varepsilon_ {u}}{d (u)} \in \mathrm{End} (\mathrm{E} _ {u}) \subset \mathrm{F} (\hat {\mathrm{G}}).\tag{25}
\]

命题 5. 设 f 是 G 上的连续中心函数。当且仅当函数 \(u \mapsto |\langle \chi_{u}|f\rangle|\) 在 \(\hat{G}\) 上速降时，f 才是无穷次可微的；在此情形下

\[
f (g) = \sum_ {u \in \hat {\mathrm{G}}} \langle \chi_ {u} | f \rangle \chi_ {u} (g)
\]

对于所有 g ∈ G。

由定理 1 b)，函数 \(\overline{f}\) 无穷次可微，当且仅当函数 \(u \mapsto \| u(\overline{f})\|_{\infty}\) 速降；但由 (20)，

\[
\left\| u (\overline {{f}}) \right\| _ {\infty} = \frac {1}{d (u)} | \langle \chi_ {u} | f \rangle |,
\]

由于函数 \(d(u)\) 和 \(\frac{1}{d(u)}\) 都是缓增的，故得第一条断言。

设 \(f\) 无穷次可微；由定理 1 a)，有 \(f(g) = \sum_{u\in \hat{\mathbf{G}}}f_u(g)\) 对所有 \(g\in \mathrm{G}\) 成立，于是

\[
\begin{array}{c} f _ {u} (g) = \langle u (g) | u (f) \rangle = d (u) \mathrm{Tr} (u (g) ^ {- 1}. u (f)) = d (u) \mathrm{Tr} \left(u (g) ^ {- 1} \langle \overline {{\chi}} _ {u} | f \rangle \frac {\varepsilon_ {u}}{d (u)}\right) \\ = \langle \overline {{\chi}} _ {u} | f \rangle \mathrm{Tr} (u (g) ^ {- 1}) = \langle \overline {{\chi}} _ {u} | f \rangle \overline {{\chi}} _ {u} (g). \end{array}
\]

因此，\(f(g) = \sum_{u\in \hat{\mathbf{G}}}\langle \overline{\chi}_u|f\rangle \overline{\chi}_u(g)\)；但对于所有 \(u\in \hat{\mathbf{G}}\)，\(u'\) 是 \(u\) 的反步表示，满足 \(\overline{\chi}_u = \chi_{u'}\)，且映射 \(u\mapsto u'\) 是 \(\hat{\mathbf{G}}\) 的一个置换；所以也有 \(f(g) = \sum_{u\in \hat{\mathbf{G}}}\langle \chi_u|f\rangle \chi_u(g)\)，命题得证。

推论。设 \(f\) 是 \(G\) 上的连续中心函数。则 \(f\) 是无穷次可微的，当且仅当 \(f\) 在 \(T\) 上的限制是无穷次可微的。

事实上，由 §7, no. 4 的 Cor. 4，

\[
\langle \chi_ {u} | f \rangle = \int_ {\mathrm{G}} \overline {{\lambda (u) (t)}} \varphi (t) d t, \quad \text { where } \varphi (t) = \prod_ {\alpha > 0} (1 - \alpha (t) ^ {- 1}) f (t).
\]

若 \(f|_{\mathrm{T}}\) 是无穷次可微的，则 \(\varphi\) 也是如此；将命题 5 应用于群 \(\mathrm{T}\)，可知函数 \(\mu \mapsto \int_{\mathrm{T}} \overline{\mu(t)} \varphi(t) dt\) 在 \(\hat{\mathrm{T}} = \mathrm{X}(\mathrm{T})\) 上速降，函数 \(u \mapsto \langle \chi_u | f \rangle\) 也速降；因此函数 \(f\) 是无穷次可微的（命题 5）。反之显然。

